%============================================================================
% OPERATIONAL EXECUTIVE SUMMARY of the Technical Reference
% Sustainable University IT --- A Technical Reference for Architecture and Migration
%
% Audience: CISO, CTO, infrastructure / network team leads (operational).
% Derived from: ../solution-EN.tex and its chapters (the "Technical Reference").
% Methodology: /companion, applied with an OPERATIONAL framing.
% Fully generic: no institution named.
%============================================================================
\documentclass[11pt,a4paper,twoside]{article}

% --- Encoding and language ---
\usepackage[utf8]{inputenc}
\usepackage[T1]{fontenc}
\usepackage[english]{babel}

% --- Page layout (compact for 3-6 pages) ---
\usepackage[top=1.8cm,bottom=1.6cm,left=2.1cm,right=2.1cm]{geometry}
\usepackage{setspace}
\setstretch{1.04}

% --- Typography ---
\usepackage{lmodern}
\usepackage{microtype}

% --- Graphics and colours ---
\usepackage{graphicx}
\usepackage{suit-logo}
\usepackage[dvipsnames,table]{xcolor}
\definecolor{Slate}{RGB}{70,70,90}
\definecolor{Crimson}{RGB}{220,20,60}
\usepackage{tikz}
\usetikzlibrary{shapes,arrows.meta,positioning,fit,calc,backgrounds}

% --- Tables ---
\usepackage{booktabs}
\usepackage{tabularx}
\usepackage{array}
\usepackage{multirow}

% --- Lists ---
\usepackage{enumitem}
\usepackage{pifont}
\setlist{nosep,leftmargin=1.5em}

% --- Links ---
\usepackage[colorlinks=true,linkcolor=NavyBlue,citecolor=OliveGreen,
            urlcolor=BrickRed]{hyperref}
\usepackage{xurl}

% --- Quotation handling (required by biblatex) ---
\usepackage{csquotes}

% --- Bibliography: biblatex + biber (shared consolidated bib) ---
\usepackage[backend=biber,
            style=numeric,
            sorting=none,
            sortcites=true,
            maxbibnames=99]{biblatex}
\addbibresource{../../shared/suit-consolidated.bib}

% --- Bibliography display modes (standard / annotated / full) ---
\usepackage{bib-modes}
\setbibmode{linked}

% --- Glossary with tooltips ---
\usepackage{gls-modes}
% ============================================================================
% Glossary definitions for gls-modes.sty
% \glsdefine{key}{SHORT}{Long Form}{Tooltip definition}
% ============================================================================
%
% SUIT NOTE: this is the institution-neutral SUIT glossary. All entries are
% role-based / generic. Adopting institutions supply their own local values
% via the Localization Guide (shared/localization-guide-EN.tex) rather than
% via hard-coded glossary entries.
% ============================================================================

% --- Technical Infrastructure ---
\glsdefine{acl}{ACL}{Access Control List}{Access Control List -- a set of rules on a network device that permits or denies traffic based on source/destination addresses and ports.}
\glsdefine{api}{API}{Application Programming Interface}{Application Programming Interface -- a standardised interface allowing software systems to communicate with each other.}
\glsdefine{byod}{BYOD}{Bring Your Own Device}{Bring Your Own Device -- the practice of employees or students using personal devices to access institutional resources.}
\glsdefine{dhcp}{DHCP}{Dynamic Host Configuration Protocol}{Dynamic Host Configuration Protocol -- automatically assigns IP addresses to devices on a network.}
\glsdefine{dmz}{DMZ}{Demilitarised Zone}{Demilitarised Zone -- a network segment between an internal network and the Internet, hosting externally accessible services.}
\glsdefine{dns}{DNS}{Domain Name System}{Domain Name System -- translates human-readable domain names into IP addresses.}
\glsdefine{dpi}{DPI}{Deep Packet Inspection}{Deep Packet Inspection -- a method of examining the content of network packets as they pass through a checkpoint.}
\glsdefine{dtn}{DTN}{Data Transfer Node}{Data Transfer Node -- a server optimised for high-throughput scientific data transfers in a Science DMZ.}
\glsdefine{edr}{EDR}{Endpoint Detection and Response}{Endpoint Detection and Response -- security software monitoring endpoints for threats.}
\glsdefine{hpc}{HPC}{High-Performance Computing}{High-Performance Computing -- computing systems designed for large-scale scientific simulations and data processing.}
\glsdefine{iam}{IAM}{Identity and Access Management}{Identity and Access Management -- systems and processes for managing digital identities and controlling access to resources.}
\glsdefine{nac}{NAC}{Network Access Control}{Network Access Control -- system enforcing authentication, device posture assessment, and dynamic VLAN assignment at the network edge (e.g. PacketFence, Cisco ISE).}
\glsdefine{ids}{IDS}{Intrusion Detection System}{Intrusion Detection System -- monitors network traffic for suspicious activity (detection only, no inline blocking).}
\glsdefine{ips}{IPS}{Intrusion Prevention System}{Intrusion Prevention System -- monitors and actively blocks detected threats inline.}
\glsdefine{lms}{LMS}{Learning Management System}{Learning Management System -- platform for delivering and managing educational content (e.g. Moodle).}
\glsdefine{mfa}{MFA}{Multi-Factor Authentication}{Multi-Factor Authentication -- requiring two or more verification factors to access a resource.}
\glsdefine{pki}{PKI}{Public Key Infrastructure}{Public Key Infrastructure -- framework for managing digital certificates and encryption keys.}
\glsdefine{rbac}{RBAC}{Role-Based Access Control}{Role-Based Access Control -- access permissions assigned to roles rather than to individual users.}
\glsdefine{sciencedmz}{Science DMZ}{Science Demilitarised Zone}{Science DMZ -- a network architecture pattern developed by ESnet for high-performance scientific data transfers, bypassing enterprise firewalls in favour of ACLs and dedicated IDS.}
\glsdefine{sdn}{SDN}{Software-Defined Networking}{Software-Defined Networking -- centralised, programmable control of network behaviour.}
\glsdefine{siem}{SIEM}{Security Information and Event Management}{Security Information and Event Management -- centralised platform for collecting and analysing security logs and alerts.}
\glsdefine{soc}{SOC}{Security Operations Centre}{Security Operations Centre -- team and facility responsible for monitoring and responding to security incidents.}
\glsdefine{sso}{SSO}{Single Sign-On}{Single Sign-On -- authentication allowing access to multiple systems with one set of credentials.}
\glsdefine{tac}{TAC}{Trusted Activity Context}{Trusted Activity Context -- a coherent, validated environment assembling data, tools, services, and access rights for a specific university activity (concept introduced in this report).}
\glsdefine{vdi}{VDI}{Virtual Desktop Infrastructure}{Virtual Desktop Infrastructure -- technology delivering desktop environments from a centralised server to any device.}
\glsdefine{vlan}{VLAN}{Virtual Local Area Network}{Virtual Local Area Network -- a logical subdivision of a physical network, isolating traffic between groups of devices.}
\glsdefine{vpn}{VPN}{Virtual Private Network}{Virtual Private Network -- encrypted tunnel for secure remote connectivity.}

% --- Governance Frameworks and Standards ---
\glsdefine{abac}{ABAC}{Attribute-Based Access Control}{Attribute-Based Access Control -- access decisions based on attributes of the user, resource, and environment (NIST SP 800-162).}
\glsdefine{cobit}{COBIT}{Control Objectives for Information and Related Technologies}{COBIT -- Control Objectives for Information and Related Technologies, IT governance framework by ISACA distinguishing governance from management.}
\glsdefine{fair}{FAIR}{Factor Analysis of Information Risk}{FAIR -- Factor Analysis of Information Risk, a quantitative methodology for assessing and managing information risk.}
\glsdefine{isms}{ISMS}{Information Security Management System}{Information Security Management System -- the set of policies and procedures for systematically managing security (ISO 27001).}
\glsdefine{zt}{Zero Trust}{Zero Trust architecture}{Zero Trust -- security model based on never trust, always verify; access decisions based on identity and device posture rather than network location (NIST SP 800-207).}

% --- European Regulations and Legal Instruments ---
\glsdefine{gdpr}{GDPR}{General Data Protection Regulation}{General Data Protection Regulation (EU) 2016/679 -- EU regulation on the protection of personal data.}
\glsdefine{nis2}{NIS2}{NIS2 Directive}{NIS2 Directive (EU) 2022/2555 -- directive on measures for a high common level of cybersecurity.}
\glsdefine{aiact}{AI Act}{European AI Act}{European AI Act -- Regulation (EU) 2024/1689 on artificial intelligence, risk-based classification of AI systems.}
\glsdefine{cjeu}{CJEU}{Court of Justice of the European Union}{Court of Justice of the European Union -- the EU highest court, interpreting EU law.}
\glsdefine{ecthr}{ECtHR}{European Court of Human Rights}{European Court of Human Rights -- court enforcing the ECHR, based in Strasbourg.}
\glsdefine{edpb}{EDPB}{European Data Protection Board}{European Data Protection Board -- EU body ensuring consistent application of the GDPR.}
\glsdefine{edps}{EDPS}{European Data Protection Supervisor}{European Data Protection Supervisor -- independent EU authority overseeing the application of GDPR by EU institutions and bodies (distinct from the EDPB which coordinates national DPAs).}
\glsdefine{enisa}{ENISA}{European Union Agency for Cybersecurity}{ENISA -- European Union Agency for Cybersecurity.}
\glsdefine{dpo}{DPO}{Data Protection Officer}{Data Protection Officer -- person responsible for overseeing GDPR compliance within an organisation.}
\glsdefine{icradp}{ICRADP}{Independent Commission for Reasonable Adjustments of Digital Policies}{ICRADP -- Independent Commission for Reasonable Adjustments of Digital Policies (proposed in this report; French: CARPN).}

% --- Cross-border data transfers ---
\glsdefine{tia}{TIA}{Transfer Impact Assessment}{Transfer Impact Assessment -- documented assessment of whether a third country offers a level of personal-data protection essentially equivalent to EU law, required by CJEU Schrems II for transfers outside the EEA when relying on Standard Contractual Clauses.}
\glsdefine{scc}{SCC}{Standard Contractual Clauses}{Standard Contractual Clauses -- pre-approved European Commission contractual clauses for transfers of personal data outside the EEA; valid under Schrems II only when supplemented by a Transfer Impact Assessment and supplementary measures.}

% --- Sovereignty grid and technical reference vocabulary ---
\glsdefine{it}{IT}{Information Technology}{Information Technology -- the systems, software, networks and processes used to store, manipulate and transmit institutional information.}
\glsdefine{scim}{SCIM}{System for Cross-domain Identity Management}{System for Cross-domain Identity Management -- IETF standard (RFC 7643/7644) for automating the exchange of user identity information between identity domains.}
\glsdefine{byok}{BYOK}{Bring Your Own Key}{Bring Your Own Key -- a key-management arrangement in which the customer generates and provides the encryption keys used by a third-party service, retaining custody and rotation control.}
\glsdefine{hyok}{HYOK}{Hold Your Own Key}{Hold Your Own Key -- an arrangement stronger than BYOK in which the customer retains the keys in their own infrastructure (typically an HSM) and the service never has cleartext access; cryptographic operations require live customer participation.}
\glsdefine{rfp}{RFP}{Request for Proposal}{Request for Proposal -- a procurement document issued by an institution to solicit proposals from suppliers, specifying functional, technical, and contractual requirements.}
\glsdefine{iot}{IoT}{Internet of Things}{Internet of Things -- networked devices typically embedded in physical objects (sensors, controllers, lab instruments) with limited or no user interface and constrained security update mechanisms.}
\glsdefine{tls}{TLS}{Transport Layer Security}{Transport Layer Security -- IETF cryptographic protocol (RFC 8446 for TLS 1.3) providing confidentiality and integrity for network communications; the foundation of HTTPS.}
\glsdefine{opa}{OPA}{Open Policy Agent}{Open Policy Agent -- CNCF graduated open-source policy decision engine using the Rego language; widely adopted as a general-purpose authorisation and admission-control PDP.}
\glsdefine{pdp}{PDP}{Policy Decision Point}{Policy Decision Point -- the component that evaluates a request against the policy corpus and returns a decision; distinct from the Policy Enforcement Point (PEP) that applies the decision.}
\glsdefine{pep}{PEP}{Policy Enforcement Point}{Policy Enforcement Point -- the component that intercepts an access request, queries the PDP, and applies the returned decision (allow/deny plus obligations).}

% --- Research Networks ---
\glsdefine{csirt}{CSIRT}{Computer Security Incident Response Team}{Computer Security Incident Response Team -- team responsible for detecting, analysing, and responding to cybersecurity incidents.}
\glsdefine{eduroam}{eduroam}{Education Roaming}{eduroam -- Education Roaming, federated Wi-Fi service allowing users to connect at any participating institution worldwide using their home credentials.}
\glsdefine{edugain}{eduGAIN}{eduGAIN interfederation service}{eduGAIN -- global interfederation service connecting identity federations, enabling cross-institutional single sign-on.}
\glsdefine{geant}{GEANT}{GEANT pan-European research network}{GEANT -- pan-European backbone network connecting 40+ NRENs, serving 50 million users in 10,000+ institutions.}
\glsdefine{nren}{NREN}{National Research and Education Network}{National Research and Education Network -- a specialised ISP dedicated to supporting research and education (each country operates one).}
\glsdefine{first}{FIRST}{Forum of Incident Response and Security Teams}{FIRST -- Forum of Incident Response and Security Teams, global network of CSIRTs facilitating coordination across national and sectoral incident-response teams.}


% === Localization (key/value layer 0: generic defaults) ===
\usepackage{localization}
% ============================================================================
% localization-defaults.tex — LAYER 0 default bindings
% ----------------------------------------------------------------------------
% One \setloc per localization variable. This file is the authoritative registry
% of localization variables (49 distinct keys — INSTITUTION_NAME and
% ACADEMIC_GOVERNING_BODIES each support two anchor families and bind once here).
%
% Each default value is the variable's CANONICAL GENERIC ROLE PHRASE: the
% jurisdiction-neutral, institution-agnostic wording already used in the
% suite's current generic prose (the genericDecision of the matching row).
% Rendering a document with ONLY these defaults loaded therefore reproduces
% the present generic text UNCHANGED. An edition file (e.g. an EU or a
% Luxembourg edition) overrides any subset of these keys with \setloc after
% this file is \input.
%
% Keys are kebab-case. Requires localization.sty (\setloc / \loc).
% ============================================================================

% --- Subject institution & profile ------------------------------------------
% edition-name is the edition identity rendered prominently on every title
% page (\suiteditionline) and in \suitmetablock. The Layer-0 value makes an
% uninstantiated build self-identify loudly rather than pass for an edition.
\setloc{edition-name}{Generic reference edition (not instantiated)}
\setloc{institution-name}{the institution}
\setloc{institution-type}{a typical research-intensive institution}
\setloc{institution-macro-set}{the institution}
\setloc{staffing-profile}{an institution that sets excellence as a strategic objective}

% --- Institutional units & pilots -------------------------------------------
\setloc{cs-research-unit}{a computer-science department}
\setloc{security-research-centre}{a cybersecurity-focused interdisciplinary research centre}
\setloc{institution-hpc-facility}{the institution's HPC facility / HPC team}
\setloc{pilot-entities}{the institution's HPC facility / team}

% --- Institutional governance roles -----------------------------------------
\setloc{executive-leadership}{the institution's central executive leadership}
\setloc{governing-board}{the institution's governing board}
\setloc{academic-governing-bodies}{the academic deliberative councils}
\setloc{academic-endorsement-path}{the relevant academic deliberative councils}
\setloc{institution-faculty-structure}{typically organised into faculties, each with departments of teaching and research, and interdisciplinary centres}
\setloc{staff-representation-body}{the institution's staff-representation bodies}
% Approving academic units for the internal-approval byline on the dossier and
% the rectorate summary (the faculty/department domains that signed off the
% institution's edition). Generic default is the neutral role phrase.
\setloc{approving-academic-units}{the approving academic units of \loc{institution-name}}

% Suite version (generic, maintainer-managed; bumped by tools/bump-version.sh).
% ============================================================================
% suite-version.tex --- SINGLE SOURCE OF TRUTH for the generic SUIT suite
% version (maintainer-managed). Common to all four generic documents.
% ----------------------------------------------------------------------------
% Semantic version X.Y.Z. Bumped ONLY by tools/bump-version.sh:
%   patch (default) -> Z+1   (one bump per release, regardless of change count)
%   minor           -> Y+1, Z=0   (maintainer, on request)
%   major           -> X+1, Y=0, Z=0   (maintainer, on request)
% Do not edit by hand other than through the bump tool; the line below is
% matched verbatim by the tool's parser.
% ============================================================================
\setloc{suite-version}{1.0.0}

% Per-document edition versions: EMPTY in the generic build (the suite version
% is shown instead); an edition fills, in its set.tex, the documents it
% publishes. Keyed by document id (derived from \jobname in \suitmetablock).
\setloc{policy-EN-edition-version}{}
\setloc{policy-summary-EN-edition-version}{}
\setloc{solution-EN-edition-version}{}
\setloc{solution-summary-EN-edition-version}{}

% Document metadata for the unified \suitmetablock title-page block. Each
% edition overrides these for its own visibility/approval facts.
\setloc{doc-visibility}{INTERNAL}
\setloc{doc-approval-authority}{\loc{approving-academic-units}}
\setloc{doc-approval-date}{\textit{(approval date to be set)}}

% --- Research & education network and its services --------------------------
\setloc{nren-name}{the national research and education network}
\setloc{nren-csirt}{the NREN CSIRT}
\setloc{national-private-sector-csirt}{a separate CSIRT serving the private sector, communes and non-governmental entities}
\setloc{financial-sector-supervisor}{a financial-sector supervisor}
\setloc{research-ecosystem-partners}{a handful of national research institutes}
\setloc{sdmz-status-lead}{In many national ecosystems, the}

% --- Companion openness-status table: one marker per element ('\warn' = to be
%     assessed in the generic build; an edition binds the institution's verdict) ---
\setloc{status-table-subtitle}{current status}
\setloc{status-byod}{\warn}
\setloc{status-eduroam}{\warn}
\setloc{status-mobility}{\warn}
\setloc{status-internet}{\warn}
\setloc{status-cloud-ai}{\warn}
\setloc{status-sci-sources}{\warn}
\setloc{status-local-admin}{\warn}
\setloc{status-tool-choice}{\warn}
\setloc{status-environments}{\warn}
\setloc{status-data-platforms}{\warn}
\setloc{status-code-repos}{\warn}
\setloc{status-open-access}{\warn}
\setloc{status-labs}{\warn}
\setloc{status-cyber-ranges}{\warn}
\setloc{status-api-access}{\warn}

% --- Abstract genesis + ICRADP adjustments-committee precedent ---
\setloc{abstract-genesis-institution}{a research-intensive university}
\setloc{existing-adjustments-committee-precedent}{Many institutions already operate a \keyword{reasonable-adjustments committee} (or equivalent)}
\setloc{icradp-committee-precedent-lead}{Where the institution already operates a reasonable-adjustments mechanism (supporting students and staff), the \gls{icradp} extends that principle to digital policies; where none exists, it establishes the principle directly.}
\setloc{cite-adjustments-committee}{}
\setloc{cite-equal-treatment-law}{}
\setloc{funding-instrument}{multi-year planning and funding cycle}
\setloc{peer-nren-roster}{peer NRENs in comparable jurisdictions, each operating a national backbone, a CSIRT, and federated identity/storage/IaaS services}
\setloc{nren-service-catalogue}{the NREN-provided federated-service catalogue (federated Wi-Fi, trusted certificate service, performance monitoring, managed research-data transfer)}
\setloc{nren-research-cloud}{a national/regional research cloud or NREN-provided cloud service, where one exists}
\setloc{research-backbone-scale}{the regional research backbone (membership and backbone capacity per region)}
\setloc{identity-federation-profile}{the academic identity interfederation}

% --- High-performance computing ---------------------------------------------
\setloc{national-hpc-facility}{a national/EuroHPC supercomputer}

% --- Deployment topology & data custody -------------------------------------
\setloc{deployment-custody-model}{the institution operates, or procures under equivalent custody guarantees, a deployment of its core components}
\setloc{data-custody-sovereignty-profile}{a custody arrangement that minimises the applicable disclosure-regime exposure}

% --- Regulatory regimes (supranational anchors as default profile) ----------
\setloc{regulatory-regime}{binding law in the EU and Council of Europe member states (including the institution's jurisdiction)}
\setloc{cybersecurity-baseline-regime}{the applicable supranational and national legal framework}
\setloc{data-protection-regime}{the institution's applicable data-protection regime}
\setloc{eid-trust-services-regime}{the institution's applicable trust-services regime}
\setloc{ai-regulation}{the adopter's applicable AI-governance regime}
\setloc{cross-border-transfer-regime}{the applicable cross-border data-transfer regime}
\setloc{essential-entity-classification-basis}{, the public-administration entry of NIS2 Annex~I being the basis in the worked example}
\setloc{nis2-sanction-ceiling}{up to EUR\,10\,M or 2\,\% of annual turnover}
% --- NIS2 essential-entity status (conditional-fact pattern; an edition asserts the established fact via its set) ---
\setloc{essential-entity-modality}{may be classified}
\setloc{nis2-status-conditional-tail}{, depending on the national transposition}
\setloc{cer-designation-status}{A given institution may or may not be designated as a CER critical entity}
% --- Edition-routed citations: empty in the generic build; bound to \cite{} markers in an edition set ---
\setloc{cite-ilr-essential}{}
\setloc{cite-nis2-transposition}{}
\setloc{cite-notification-platform}{}
\setloc{cite-national-risk-method}{}
\setloc{cite-public-sector-isp}{}
\setloc{cite-cyber-strategy}{}
\setloc{cite-nren-csirt}{}
\setloc{cite-cer-transposition}{}
\setloc{cite-gov-csirt}{}
\setloc{cite-spoc}{}
\setloc{cite-infosec-agency}{}
\setloc{cite-circl}{}
\setloc{cite-cssf}{}
\setloc{cite-restena}{}

% --- National transposition instruments -------------------------------------
\setloc{nis2-national-transposition}{the national NIS2 transposition law}
\setloc{cer-national-transposition}{the national CER transposition law}
\setloc{national-equal-treatment-law}{the national transposition of Directive 2000/78/EC into labour law}
\setloc{national-public-sector-isp}{a national public-sector information-security policy}
\setloc{national-cyber-strategy}{the national cybersecurity strategy}
\setloc{national-risk-method}{a national risk-analysis method}
\setloc{national-cyber-commissariat-founding-law}{the national founding statute of the cyber-resilience commissariat}

% --- National authorities & bodies (role-based) -----------------------------
\setloc{national-dpa}{the competent data-protection authority}
\setloc{national-gov-csirt}{the national/governmental CSIRT}
\setloc{national-infosec-agency}{the national information-security agency}
\setloc{nis2-supervisory-authority}{the national NIS2 supervisory authority}
\setloc{national-spoc-authority}{the national single point of contact}
\setloc{national-cyber-competence-centre}{a national cybersecurity competence centre}
\setloc{national-notification-platform}{the national NIS2 incident-notification platform}
\setloc{national-legislature}{the national legislature}
\setloc{national-executive-department}{the national executive department responsible for the protection commissariat}

% --- Operational contacts & market exemplars --------------------------------
\setloc{internal-emergency-incident-contact}{the institution's internal emergency / incident-reporting contact}
\setloc{incumbent-identity-suite-vendor}{a pure-commercial identity provider}

% --- ICRADP administrative-review service levels ----------------------------
% Adjustment-commission decision deadline and emergency provisional-grant
% window. Default to the worked figures; an edition may rebind them to its own
% administrative-review service levels.
\setloc{icradp-decision-deadline}{15 working days}
\setloc{icradp-emergency-deadline}{48 hours}

% --- Assurance / review recency window (single-source magic number) ---------
% The audit/pen-test/drill recency window used by the Chapter 4 self-assessment.
% Centralised so a re-tuned interval is one edit and an edition may rebind it.
\setloc{assurance-recency-window}{twenty-four months}

\endinput


% --- Boxes ---
\usepackage{tcolorbox}
\tcbuselibrary{skins,breakable}

% --- Headers ---
\usepackage{fancyhdr}
\setlength{\headheight}{15pt}
\pagestyle{fancy}
\fancyhf{}
\fancyhead[LO,RE]{\small\textit{Operational Summary --- Sustainable University IT}}
\fancyhead[LE,RO]{\small\thepage}
\fancyfoot[C]{\suitfootlogo}
\renewcommand{\headrulewidth}{0.4pt}
\renewcommand{\footrulewidth}{0.4pt}

% --- Semantic commands (aligned with the rest of the suite) ---
\newcommand{\keyword}[1]{\textbf{#1}}
\newcommand{\institution}[1]{\textsc{#1}}
\newcommand{\framework}[1]{\textsf{#1}}

% --- No section numbering ---
\setcounter{secnumdepth}{0}

% --- Compact sections ---
\usepackage{titlesec}
\titlespacing*{\section}{0pt}{0.8em}{0.3em}
\titlespacing*{\subsection}{0pt}{0.5em}{0.2em}
\titleformat{\section}{\normalsize\bfseries\color{NavyBlue}}{}{}{}
\titleformat{\subsection}{\small\bfseries\color{MidnightBlue}}{}{}{}

%----------------------------------------------------------------------------
% Cross-reference helpers back to the full Technical Reference (SUITSOLUTION).
% Self-contained: the summary names the target chapter explicitly so that the
% pointer resolves without building the full Reference first.
%----------------------------------------------------------------------------
% \trchap{N}{title} -> "Chapter N (title) of the Technical Reference"
\newcommand{\trchap}[2]{%
  Chapter~#1%
  \ifx\\#2\\\else\space(\emph{#2})\fi\space of the Technical Reference%
}
% \mainref{...} -> grey italic pointer line into the full Reference
\newcommand{\mainref}[1]{%
  \par\vspace{1pt}
  {\footnotesize\color{gray}\itshape
  $\triangleright$~Full treatment, contracts and references: #1,
  in the Technical Reference~\cite{ng-tech-ref-2026}.}\par\vspace{2pt}}

% --- Operational box palette (consistent with the suite) ---
% principle (NavyBlue) / alert (BrickRed) / action (OliveGreen)
\newtcolorbox{principlebox}[1][]{%
  colback=NavyBlue!5!white, colframe=NavyBlue!70!black,
  fonttitle=\bfseries, title={#1}, boxrule=0.5pt, sharp corners=south,
  left=6pt, right=6pt, top=4pt, bottom=4pt}

\newtcolorbox{alertbox}[1][]{%
  colback=BrickRed!5!white, colframe=BrickRed!70!black,
  fonttitle=\bfseries, title={#1}, boxrule=0.5pt, sharp corners=south,
  left=6pt, right=6pt, top=4pt, bottom=4pt}

\newtcolorbox{actionbox}[1][]{%
  colback=OliveGreen!5!white, colframe=OliveGreen!70!black,
  fonttitle=\bfseries, title={#1}, boxrule=0.5pt, sharp corners=south,
  left=6pt, right=6pt, top=4pt, bottom=4pt}

% --- Prevent vertical stretching ---
\raggedbottom

%============================================================================
\begin{document}

% --- Title block (no separate title page --- space is precious) ---
\begin{center}
\vspace*{-0.2cm}

{\suittitlelogo[0.30\textwidth]}

\smallskip
{\Large\bfseries Sustainable University IT --- Operational Summary}

\smallskip
{\normalsize\itshape Reference Architecture, Migration Waves and Acceptance Gates,
at a Glance}

\smallskip
{\small For the CISO, the CTO, and infrastructure / network team leads}

\smallskip
{\scriptsize Operational companion to: \emph{``Sustainable University IT ---
A Technical Reference for Architecture and Migration''}}

\smallskip
\suitmetablock

\medskip

\begin{tcolorbox}[
    colback=NavyBlue!8!white, colframe=NavyBlue!85!black,
    boxrule=1.2pt, arc=3pt,
    left=10pt, right=10pt, top=8pt, bottom=8pt, halign=center]
{\normalsize\bfseries Nine layers, three planes, ten waves ---
each wave reversible, each closure proven\\[3pt]
by conformance, every component replaceable through a published interface.}
\end{tcolorbox}
\end{center}

\smallskip
\noindent\rule{\textwidth}{0.6pt}

%============================================================================
\section{Operational synopsis}
%============================================================================

\noindent We are standing up a \keyword{vendor-neutral reference
architecture} for university IT and a \keyword{brownfield migration
playbook} to reach it from whatever estate exists today. In operational
terms: a small, strongly protected core (identity, observability, network
backbone) is established first; authorisation is lifted out of scattered
application logic and identity-provider-native rules into a single
\gls{pdp}; the network is carved into declared zones; and every workload,
session and access decision is bound to a \gls{tac} --- a validated bundle of
subject, action, resource and context that the platform can activate,
audit and tear down as a unit. The architecture's normative content is a
set of \keyword{interface contracts} (OIDC/SCIM, RADIUS, NETCONF/YANG, the
Kubernetes and S3 APIs, OpenTelemetry, KMIP/ACME), not a product list:
open-source instantiations prove the contracts are realisable, and any
compliant component --- proprietary or open --- is admissible.
The migration is delivered as \keyword{ten waves} under a dependency graph,
each with a coexistence (shadow / dual-stack) pattern, a declared
reversibility window, and conformance-based exit gates. The pay-off for
operations is concrete: incidents are bounded by TAC instead of cascading
through a flat perimeter, suppliers are swapped through a rehearsed exit
rather than a forklift, and ``done'' is a passing test rather than a
declaration.

\mainref{\trchap{1}{Synthesis}}

%============================================================================
\section{Reference architecture at a glance: nine layers, three planes}
%============================================================================

\begin{table}[h!]
\centering\small
\renewcommand{\arraystretch}{1.15}
\begin{tabularx}{\textwidth}{
  >{\bfseries\raggedright\arraybackslash}p{3.0cm}
  >{\raggedright\arraybackslash}p{3.3cm}
  >{\raggedright\arraybackslash}X}
\toprule
\textbf{Plane} & \textbf{Layers} & \textbf{Operational role / primary contract} \\
\midrule
Subject \& Trust &
  Identity (Ch.~6), Endpoint (Ch.~11), Crypto (Ch.~13) &
  Who is acting, from where, with what assurance. OIDC/SAML2/SCIM~2.0;
  posture schema; PKCS\#11/ACME/KMIP. Read by the Control plane. \\
\addlinespace
Execution &
  Network (Ch.~8), Compute (Ch.~9), Storage (Ch.~10) &
  Where activity runs and on which data. RADIUS + NETCONF/YANG + IPFIX;
  Kubernetes/OpenStack/OCI; S3/NFS/KMIP. Written by the Control plane. \\
\addlinespace
Control \& Governance &
  Policy (Ch.~7), Observability (Ch.~12), Orchestration (Ch.~14) &
  How decisions are made, recorded and executed. \gls{pdp} verdict
  \texttt{\{verdict, obligations, policy\_version\}}; OpenTelemetry +
  canonical TAC schema; cross-layer TAC activation (no single standard). \\
\bottomrule
\end{tabularx}
\caption{The reference architecture: nine functional layers grouped in three
planes. Two regular flows govern it --- the Subject \& Trust plane feeds the
Control plane (identity, posture and crypto parameterise policy); the Control
plane writes to the Execution plane (decisions enforced at Network, Compute,
Storage; Orchestration realises them; Observability collects evidence of every
crossing).}
\end{table}

\begin{principlebox}[Operating rule: contracts, not products]
Procurement and acceptance reference \emph{standards}, not product names. A
component is in production only when it speaks its layer's interface contract
and is \keyword{replaceable}: standard interfaces with no captive extensions
in the critical path, open state/config formats, a documented exit-data
export within an agreed window, and no hidden supplier sidecars. The one
honest gap is Orchestration: no single standard yet exists, so the
cross-layer ``glue'' is bounded deliberately to avoid creating a new captive
dependency.
\end{principlebox}

\mainref{\trchap{5}{Architectural principles} and the per-layer
chapters~6--14}

%============================================================================
\section{Migration sequence: ten waves, dependencies, reversibility}
%============================================================================

The waves are not calendar phases; they unfold under a dependency graph that
permits parallelism. \keyword{Wave~0} is the entry point for every
trajectory. After it, Waves~1--3 (Identity, Observability, Network) may run
in parallel if staffing allows; Wave~4 (Policy) needs Identity and
Observability; Wave~5 (Pilot) is the first end-to-end exercise on a contained
scope. \keyword{Every wave is reversible at acceptable cost during a published
window} --- the coexistence pattern is the reversal mechanism: as long as the
incumbent stays viable, the wave is undone by re-routing to it.

\begin{table}[h!]
\centering\small
\renewcommand{\arraystretch}{1.12}
\begin{tabularx}{\textwidth}{
  >{\bfseries\raggedright\arraybackslash}p{2.2cm}
  >{\raggedright\arraybackslash}X
  >{\raggedright\arraybackslash}p{1.5cm}
  >{\raggedright\arraybackslash}p{3.0cm}}
\toprule
\textbf{Wave} & \textbf{Operational goal} & \textbf{Requires} &
  \textbf{Reversibility window} \\
\midrule
0 --- Foundations &
  Mandate, governance body, sovereignty grid adopted, baseline inventory.
  No technical change. & --- &
  Highest; decisions informational, not binding. \\
\addlinespace
1 --- Identity &
  Sovereign IdP on OIDC/SAML2/SCIM, R\&E federation, reviewable attribute
  schema, ingestible auth log. & 0 &
  Throughout cutover; $\sim$6--12 months after final cohort. \\
\addlinespace
2 --- Observability &
  Collector layer emitting canonical TAC schema; portable ingestion;
  retention schedule. & 0 &
  Structural; remove collectors to revert. \\
\addlinespace
3 --- Network &
  Declared zone catalogue, RADIUS attachment, NETCONF/YANG config, IPFIX,
  Science~DMZ where warranted. & 0 (rec.~2) &
  Per-segment, $\sim$3--6 months after each cut-over. \\
\addlinespace
4 --- Policy &
  \gls{pdp} deployed; selected decisions migrated from app/IdP logic;
  authoring/review process. & 1, 2 (rec.~3) &
  Per-policy, while shadow observation continues. \\
\addlinespace
5 --- Pilot &
  One or more TACs in production on contained scope; extract learnings. &
  1,2,3,4 & Structural; stop the pilot at any time. \\
\addlinespace
6 --- Endpoint &
  Unified posture schema across OSes routed into the policy plane;
  decommissioning integrated. & 1, 4 (rec.~2,5) &
  Per-policy and per-population. \\
\addlinespace
7 --- Audit \& IR &
  TAC-correlated investigation, refreshed IR runbooks, routine audit
  replay. & 2,4,5 & Structurally additive; revert at boundary. \\
\addlinespace
8 --- Workload &
  Classification as authoritative input to provisioning, storage binding,
  policy; cluster-separation rule. & 1,3,4 (rec.~5) &
  Per-workload, for the coexistence duration. \\
\addlinespace
9 --- Lifecycle &
  Self-service TAC provisioning, GitOps as source of truth, stability
  metrics. & 1--8 in some form &
  Per-consumer-group. \\
\bottomrule
\end{tabularx}
\caption{Wave sequence with required pre-requisites and indicative
reversibility windows. Pace options: \emph{5-year compressed}
(parallelised, higher risk tolerance), \emph{7-year balanced} (recommended
default), \emph{10-year gradual} (renewal-aligned, lower risk). Trajectory
varies by starting estate (greenfield, brownfield single-suite-centric,
brownfield hybrid, selective, deferred).}
\end{table}

\begin{principlebox}[Minimal coherent core --- if resources are constrained]
A constrained team may legitimately adopt \keyword{less}, not later: Wave~0
(governance/inventory) + Wave~1 (sovereign identity) + Wave~2 (observability
/ audit trail), with Wave~4 (policy) the natural next step. This delivers
exactly the size-invariant, legally compelled functions (authentication,
audit, segmentation discipline under the applicable data-protection
regime --- \gls{gdpr} in the EU, with \gls{nis2} where in scope) and is a legitimate \emph{steady state}, not an abandoned
migration. Defer or mutualise the rest along the same dependency graph.
\end{principlebox}

\mainref{\trchap{17}{How to read this playbook} and the wave
chapters~18--27}

%============================================================================
\section{Conformance \& replaceability: the acceptance gates}
%============================================================================

A contract that cannot be checked is a slogan. Each wave closes by passing
the conformance subset for its scope --- never by the judgement that ``the
migration is done''. This is the single mechanism that prevents the classic
failure mode: \keyword{a product swap disguised as architectural progress}.

\begin{table}[h!]
\centering\small
\begin{tabularx}{\textwidth}{
  >{\bfseries\raggedright\arraybackslash}p{3.0cm} X}
\toprule
\textbf{Test category} & \textbf{What the gate verifies (operationally)} \\
\midrule
Interface fidelity (F) &
  The component speaks the standard correctly (e.g.\ OIDC discovery metadata
  conformant; Kubernetes upstream conformance passes; NETCONF apply without
  the vendor console). \\
\addlinespace
Data portability (P) &
  Data/config export in a documented open format and re-import into an
  alternative without loss (SCIM export to an alternative IdP; OCI image
  digest preserved; KMIP key migration). \\
\addlinespace
Exit rehearsal (X) &
  The replacement is \emph{rehearsed}, not simulated: an alternative restores
  service on a representative subset within the documented window. \\
\addlinespace
Decision auditability (A) &
  A recorded decision replays against its recorded inputs and policy version
  to the identical verdict; TAC deactivation reverses every recorded change. \\
\bottomrule
\end{tabularx}
\caption{The four acceptance-gate categories, mapped to the four conditions
of replaceability. Wave-to-test mapping is published per wave; e.g.\ Wave~1
closes against the four Identity tests, Wave~3 against the four Network
tests, Wave~5 against the union exercised by the pilot.}
\end{table}

\begin{principlebox}[How to run it]
Run a time-boxed \keyword{conformance day}: one engineer per affected layer,
one observability operator to retrieve replay material, a note-taker, the
scope owner. Half a day per layer; one day for a full review. Output is a
structured report (pass / conditional pass / fail, with evidence and
remediation owner) and an \keyword{exit-confidence index} per layer on the
0--3 grid scale. The index is a \emph{maturity trajectory}, not a day-one
gate: a layer consumed as a managed/mutualised service answers
replaceability through the provider's contractual exit terms, not a
self-run test.
\end{principlebox}

\mainref{\trchap{16}{Conformance and replaceability tests}}

%============================================================================
\section{Sizing reference points (research-intensive reference profile) for capacity planning}
%============================================================================

\begin{principlebox}[Status of these figures]
Indicative orders of magnitude for a \keyword{research-intensive}
reference profile: $\sim$25{,}000 active subjects, $\sim$1{,}500 elevated-access
staff, $\sim$500 active research projects. These are \emph{planning aids,
not measurements}; size your own estate by direct workload measurement.
\end{principlebox}

\begin{table}[h!]
\centering\small
\begin{tabularx}{\textwidth}{
  >{\bfseries\raggedright\arraybackslash}p{3.0cm}
  >{\raggedright\arraybackslash}X
  >{\raggedright\arraybackslash}p{2.6cm}}
\toprule
\textbf{Layer} & \textbf{Indicative reference points} & \textbf{Run FTE
(shared)} \\
\midrule
Identity &
  Auths/s: a few--tens steady, hundreds at peak; concurrent sessions
  thousands; \loc{identity-federation-profile} tens--hundreds of thousands/day. & 2--4 \\
\addlinespace
Policy &
  Tens--low hundreds decisions/s; $<$10\,ms p95 online decision; bundle
  refresh in minutes. & 1--2 \\
\addlinespace
Network &
  Backbone 40--400\,Gb/s aggregate; Science~DMZ at backbone order;
  IPFIX hundreds of millions of records/day. & 3--5 \\
\addlinespace
Compute / Storage &
  OpenStack + Kubernetes; Ceph multi-PB; per-class cluster separation. &
  4--7 / 2--4 \\
\addlinespace
Observability &
  Tens--hundreds of millions of events/day; tens--low hundreds of GB/day
  compressed; 90d hot / 1y warm / multi-year cold. & 3--5 \\
\bottomrule
\end{tabularx}
\caption{Selected sizing reference points and indicative run-time
engineer-equivalents per layer (ranges are \emph{not additive}; staff are
shared across layers). The build phase adds $\sim$1--2 FTE per active wave.}
\end{table}

\noindent\textbf{Scaling breakpoints} (reshape the instantiation, not the
architecture): Identity $>$\,$\sim$50k subjects $\to$ horizontal
partitioning; Audit $>$\,$\sim$500\,GB/day $\to$ hot/warm/cold tier
separation; Policy $>$\,$\sim$5{,}000 decisions/s $\to$ per-app sidecars;
Storage $>$ multi-PB $\to$ Ceph federation. \textbf{Down-scaling} reads the
same logic backward: a teaching-led institution collapses planes to logical
segments on shared hardware, mutualises compute/storage via a national
research cloud or R\&E service, and defers the optional Science~DMZ ---
a minimal-core envelope of $\sim$3--5 shared engineer-equivalents.

\mainref{\trchap{15}{Sizing reference points}}

%============================================================================
\section{Top operational risks, mitigations, and coexistence patterns}
%============================================================================

\begin{table}[h!]
\centering\small
\newcommand{\warn}{\textcolor{BrickRed}{\ding{72}}\,}%
\begin{tabularx}{\textwidth}{
  >{\raggedright\arraybackslash}p{4.4cm} X}
\toprule
\textbf{Operational risk} & \textbf{Mitigation} \\
\midrule
\warn Product swap mistaken for architectural progress &
  Close every wave on its conformance subset (F/P/X/A); a new IdP that fails
  data-portability and exit-rehearsal has not closed the wave. \\
\addlinespace
\warn Long-running dual stack with no exit date &
  Declare the exit date at wave open; track it in programme governance;
  escalate on drift. \\
\addlinespace
\warn Bridge fails silently (sync job stops, metadata expires) &
  Name the bridge owner explicitly at wave open, with a documented
  escalation path; surface discrepancies through observability. \\
\addlinespace
\warn Bidirectional sync without conflict resolution &
  Prefer unidirectional sync; where bidirectional is unavoidable, document
  and test an explicit conflict-resolution rule before production. \\
\addlinespace
\warn Coexistence frozen as de-facto architecture &
  Periodic review under the sovereignty grid; a coexistence unrevisited for
  24 months is re-evaluated against current cost/benefit. \\
\addlinespace
\warn Reversibility forfeited under delivery pressure &
  Forfeit only on explicit written acceptance by the competent governing
  authority, against the grid, revisited at every renewal. \\
\bottomrule
\end{tabularx}
\caption{Top operational risks during migration and their mitigations.}
\end{table}

\begin{principlebox}[Shadow-mode and coexistence patterns to standardise on]
\begin{itemize}[nosep,leftmargin=1.2em]
  \item \textbf{Incumbent directory $\leftrightarrow$ institutional IdP:}
    directory stays authoritative; new IdP federates via SAML2/OIDC;
    \emph{unidirectional} sync only.
  \item \textbf{Policy shadow mode:} the \gls{pdp} logs decisions with the
    same fidelity as enforcing mode before any decision becomes binding ---
    the safe path to flip enforcement.
  \item \textbf{Per-volume \gls{siem} $\leftrightarrow$ open store:} hot
    investigative window in the commercial \gls{siem}; long-tail retention in
    the open store; dashboards migrate on financial, not technical, triggers.
  \item \textbf{Commercial hypervisor $\leftrightarrow$ open substrate;
    enterprise array $\leftrightarrow$ distributed storage:} new workloads
    default to the open side; legacy migrates at its natural lifecycle
    boundary, validated by checksum.
  \item \textbf{eduroam $\leftrightarrow$ TAC-aware Wi-Fi:} keep the federated-roaming
    profile (eduroam) inviolate; add institutional/TAC attributes only on
    the non-roaming path.
\end{itemize}
\end{principlebox}

\mainref{\trchap{29}{Reversibility playbook} and
\trchap{30}{Coexistence patterns library}}

%============================================================================
\section{What this means for the operational teams}
%============================================================================

\begin{actionbox}[Operational call to action]
\begin{enumerate}[leftmargin=1.5em]
  \item \textbf{Start with Wave~0:} adopt the sovereignty grid as the
    procurement/acceptance instrument and complete the baseline inventory ---
    organisational, zero technical risk, but every later wave inherits it.
  \item \textbf{Stand up the minimal core} (Identity + Observability), then
    the Policy plane in shadow mode, before touching the network fabric.
  \item \textbf{Make conformance the definition of done:} schedule
    conformance days, publish the exit-confidence index, and refuse to close
    a wave on declaration alone.
  \item \textbf{Write procurement against contracts, not products,} and
    rehearse the exit at least once per layer per year --- replaceability is
    only real once exercised.
\end{enumerate}
\end{actionbox}

\vspace{0.4em}
\noindent\rule{\textwidth}{0.4pt}

{\footnotesize\textbf{Full document:} \emph{``Sustainable University IT --- A
Technical Reference for Architecture and Migration''}~\cite{ng-tech-ref-2026}
develops the five-dimension sovereignty grid, the Trusted Activity Context,
the nine-layer / three-plane reference architecture on open foundations, the
ten-wave migration playbook with per-wave trajectories and reversibility, the
conformance suite, the sizing reference points, and the regulatory mapping
(in the EU instantiation: \gls{gdpr}, \gls{nis2}, the EU AI Act, eIDAS~2, Schrems~II) against the
jurisdiction-neutral standards spine. It is the shared technical baseline for
the international university network anticipated in phase~2 of the SUIT
initiative.}

\smallskip

{\footnotesize\textbf{Note on the process:}
\textit{This document is an always living document. It is provided for
constructive feedback and revision in order to improve its quality. The
engineering of this document was partially supported by digital tools, some
of which incorporate Artificial Intelligence features.}}

\vspace{0.6em}
\noindent\rule{\textwidth}{0.4pt}
\printbibliography[title={\normalsize References}]

\end{document}
